% ===================================================================
%  Αρχείο: 15391.tex (Fragment)
%  Αυτόματη μετατροπή μέσω doc2latex.py
% ===================================================================

ΘΕΜΑ 3

Στο παρακάτω σχήμα δίνονται οι γραφικές παραστάσεις των συναρτήσεων $f(x) = \eta\mu x$ και $g(x) = \sigma\upsilon\nu x$, $x \in \lbrack 0,\ 2\pi\rbrack$.

\begin{alist}
\item Να βρείτε τις συντεταγμένες των σημείων Α και Β.

\monades{10}

\item Να βρείτε την μονοτονία της συνάρτησης $g(x)$ στο $\left\lbrack \frac{\pi}{2},\pi \right\rbrack$ και την μονοτονία της συνάρτησης $f(x)$ στο $\left\lbrack \frac{3\pi}{2},2\pi \right\rbrack$.

\monades{4}

\item Με την βοήθεια του ερωτήματος β) ή με όποιον άλλο τρόπο θέλετε, να συγκρίνετε, με δικαιολόγηση, τους αριθμούς:

\end{alist}
\begin{alist}
\item
  $\sigma\upsilon\nu\left( \frac{2\pi}{3} \right)$ και $\sigma\upsilon\nu\left( \frac{5\pi}{6} \right)$.

\monades{5}


\end{alist}
\begin{alist}
\item
  $\eta\mu\left( \frac{5\pi}{3} \right)$ και $\eta\mu\left( \frac{11\pi}{6} \right)$.

\monades{6}
\end{alist}
